% !TeX root = ../教材题库_全部条目_转写工作稿.tex
% Source: XJ-B1-C04-EX-P108-Y0840
% 第4章 / 4.2 指数函数 / 教材 p.108 / PDF 第115页
% OCR未校对：公式与文字均可能有误；未完成数学TeX转写；禁止直接发给学生
% 原题及相邻原解截图（核对用）：../题目原图/XJ-B1-C04-EX-P108-Y0840-01.webp
% 本文件只在转写完成后作为题干使用；OCR 条目尚不是可靠题干。
\begin{bankitem}{4.2 例题 例题2}{第4章 / 4.2 指数函数 / 教材 p.108 / PDF 第115页}{XJ-B1-C04-EX-P108-Y0840}{OCR未校对：公式与文字均可能有误；未完成数学TeX转写；禁止直接发给学生}
\begin{ocrdraft}
[教材p.108 / PDF第115页]
\par
0二0]医学中常用的钴60射线，穿过厚度为1cm的铅板后，强度变为原来的
\par
0．568倍，穿过厚度为cm的铅板后的强度与原来的强度之比为H(x)=aX.若铅
\par
板厚度为12cm，射线穿过铅板后的强度与原来的强度之比是多少？
\par
解由H（1）一一0，568，得H（）一0．568工．
\par
故利用计算器可算得射线穿过厚度为12的铅板后的强度与原来的强度之
\par
比是
\par
H(12）=0．568120．001128，
\par
即约为原来的千分之一
\end{ocrdraft}
\end{bankitem}
